\section{Konu 3: Vektör Fonksiyonları ve Uzay Eğrileri}

Bu bölümde 09.10 ve 15.10 tarihli defter sayfalarında ve 23.10 tarihli tahta fotoğraflarında çözülen bütün sorular, derste
işlenme sırasıyla verilmiştir. Aynı soru birden fazla fotoğrafta geçiyorsa bütün kaynaklar başlıkta yazılıdır. Sonda, derste
çalışma sorusu olarak verilen problemlerin ve kaynak kitaplardaki tipik alıştırmalara benzeyen ek alıştırmaların tam çözümleri vardır.

\begin{bilgi}[Hatırlatma]
Vektör fonksiyonlarda limit, süreklilik, türev ve integral \emph{bileşen bileşen} yapılır:
$\mathbf r(t)=x(t)\mathbf i+y(t)\mathbf j+z(t)\mathbf k$ ise $\mathbf r'(t)=x'(t)\mathbf i+y'(t)\mathbf j+z'(t)\mathbf k$,
$\int\mathbf r\,dt=\big(\int x\,dt\big)\mathbf i+\big(\int y\,dt\big)\mathbf j+\big(\int z\,dt\big)\mathbf k+\mathbf C$.
\end{bilgi}

%======================================================================
\subsection{Vektör fonksiyonları: tanım kümesi, işlemler, limit, süreklilik}

\begin{soru}[Soru 1 \foto{IMG-20261009-WA0016}, \foto{IMG-20261023-WA0060}]
$\mathbf F(t)=(t+1)\,\mathbf i+\ln(1-t)\,\mathbf j+\sqrt{4-t^2}\,\mathbf k$ fonksiyonunun tanım kümesini bulunuz.
\end{soru}
\begin{cevap}[Çözüm (adım adım)]
\textbf{Fikir:} $\mathbf F(t)$'nin tanımlı olması için üç bileşenin de tanımlı olması gerekir: $T_{\mathbf F}=T_f\cap T_g\cap T_h$.

\textbf{Adım 1.} $f(t)=t+1$ polinomdur, her reel sayıda tanımlı: $T_f=(-\infty,+\infty)$.

\textbf{Adım 2.} $g(t)=\ln(1-t)$: logaritmanın argümanı pozitif olmalı. $1-t>0\iff t<1$, \ $T_g=(-\infty,1)$.

\textbf{Adım 3.} $h(t)=\sqrt{4-t^2}$: kök içi negatif olamaz. $4-t^2\ge0\iff t^2\le4\iff-2\le t\le2$, \ $T_h=[-2,2]$.

\textbf{Adım 4.} Kesişim (sayı doğrusunda üçünün ortak kısmı):
\[ T_{\mathbf F}=(-\infty,+\infty)\cap(-\infty,1)\cap[-2,2]=\boxed{[-2,1)} . \]
$-2$ dahildir ($\sqrt0=0$), $1$ dahil değildir ($\ln0$ tanımsız).
\end{cevap}

\begin{soru}[Soru 2 \foto{IMG-20261009-WA0017}, \foto{IMG-20261023-WA0060}, \foto{IMG-20261023-WA0035}]
$\mathbf F(t)=t\,\mathbf i+t^2\mathbf j+t^3\mathbf k$, $\mathbf G(t)=(1-t)\,\mathbf i+(1+t^2)\,\mathbf j+t\,\mathbf k$, $f(t)=t+1$, $g(t)=t-2$ için
$(\mathbf F+\mathbf G)(t)$, $(\mathbf F\cdot\mathbf G)(t)$, $(\mathbf F\times\mathbf G)(t)$, $(f\mathbf G)(t)$ ve $(\mathbf F\circ g)(t)$'yi bulunuz.
\end{soru}
\begin{cevap}[Çözüm (adım adım)]
\textbf{Toplam:} karşılıklı bileşenler toplanır.
\[ (\mathbf F+\mathbf G)(t)=(t+1-t)\,\mathbf i+(t^2+1+t^2)\,\mathbf j+(t^3+t)\,\mathbf k=\mathbf i+(2t^2+1)\,\mathbf j+(t^3+t)\,\mathbf k . \]
\textbf{Skaler çarpım:} $x$'ler, $y$'ler, $z$'ler çarpılıp toplanır; sonuç reel fonksiyondur.
\[ (\mathbf F\cdot\mathbf G)(t)=t(1-t)+t^2(1+t^2)+t^3\cdot t=t-t^2+t^2+t^4+t^4=2t^4+t . \]
\textbf{Vektörel çarpım:} determinant birinci satıra göre açılır.
\begin{align*}
(\mathbf F\times\mathbf G)(t)&=\begin{vmatrix}\mathbf i&\mathbf j&\mathbf k\\t&t^2&t^3\\1-t&1+t^2&t\end{vmatrix}\\
&=\mathbf i\big(t^3-t^3(1+t^2)\big)-\mathbf j\big(t^2-t^3(1-t)\big)+\mathbf k\big(t(1+t^2)-t^2(1-t)\big)\\
&=-t^5\,\mathbf i-(t^4-t^3+t^2)\,\mathbf j+(2t^3-t^2+t)\,\mathbf k .
\end{align*}
\textbf{Skalerle çarpım:}
\[ (f\mathbf G)(t)=(t+1)\big[(1-t)\mathbf i+(1+t^2)\mathbf j+t\,\mathbf k\big]=(1-t^2)\,\mathbf i+(t^3+t^2+t+1)\,\mathbf j+(t^2+t)\,\mathbf k . \]
\textbf{Bileşke:} $\mathbf F$'de $t$ yerine $g(t)=t-2$ yazılır.
\[ (\mathbf F\circ g)(t)=\mathbf F(t-2)=(t-2)\,\mathbf i+(t-2)^2\,\mathbf j+(t-2)^3\,\mathbf k . \]
\textbf{Kontrol:} $t=1$'de $\mathbf F=(1,1,1)$, $\mathbf G=(0,2,1)$, $\mathbf F\times\mathbf G=(1\cdot1-1\cdot2,\ 1\cdot0-1\cdot1,\ 1\cdot2-1\cdot0)=(-1,-1,2)$;
formülümüz de $(-1,-(1-1+1),2-1+1)=(-1,-1,2)$ verir. \checkmark
\end{cevap}

\begin{soru}[Soru 3 \foto{IMG-20261009-WA0017}, \foto{IMG-20261023-WA0035}]
$\mathbf F(t)=(t+t^2)\,\mathbf i+2t\,\mathbf j-4t\,\mathbf k$ ve $\mathbf G(t)=(2+4t)\,\mathbf i+(1-3t)\,\mathbf j+(1+t^2)\,\mathbf k$ ortogonal midir?
\end{soru}
\begin{cevap}[Çözüm (adım adım)]
Ortogonallik için \emph{her} $t$'de $\mathbf F(t)\cdot\mathbf G(t)=0$ olmalıdır.
\begin{align*}
\mathbf F\cdot\mathbf G&=(t+t^2)(2+4t)+2t(1-3t)+(-4t)(1+t^2)\\
&=(2t+4t^2+2t^2+4t^3)+(2t-6t^2)+(-4t-4t^3)\\
&=(2t+2t-4t)+(6t^2-6t^2)+(4t^3-4t^3)=0 .
\end{align*}
Özdeş olarak sıfırdır: fonksiyonlar \textbf{ortogonaldir}.
\end{cevap}

\begin{soru}[Soru 4 \foto{IMG-20261009-WA0018}, \foto{IMG-20261023-WA0035}, \foto{IMG-20261023-WA0046}]
$\mathbf F(t)=2\cos t\,\mathbf i+2\sin t\,\mathbf j+3t\,\mathbf k$ fonksiyonunun normunu bulunuz.
\end{soru}
\begin{cevap}
\[ \|\mathbf F(t)\|=\sqrt{(2\cos t)^2+(2\sin t)^2+(3t)^2}=\sqrt{4(\cos^2t+\sin^2t)+9t^2}=\sqrt{4+9t^2} . \]
\end{cevap}

\begin{soru}[Soru 5 \foto{IMG-20261009-WA0018}, \foto{IMG-20261023-WA0036}]
$\mathbf F(t)=(1-t^2)\,\mathbf i+\dfrac{\sin8t}{t}\,\mathbf j+(e^t+1)\,\mathbf k$ için $\displaystyle\lim_{t\to0}\mathbf F(t)$'yi bulunuz.
\end{soru}
\begin{cevap}[Çözüm (adım adım)]
Limit bileşen bileşen alınır.
\begin{itemize}
\item $\mathbf i$: $\lim_{t\to0}(1-t^2)=1$.
\item $\mathbf j$: $\dfrac{\sin8t}{t}=8\cdot\dfrac{\sin8t}{8t}$ ve $u=8t\to0$ iken $\dfrac{\sin u}{u}\to1$; limit $8$.
\item $\mathbf k$: $\lim_{t\to0}(e^t+1)=2$.
\end{itemize}
\[ \lim_{t\to0}\mathbf F(t)=\mathbf i+8\,\mathbf j+2\,\mathbf k . \]
\textit{Not:} Tahtada limit ``$t\to t_0$'' diye yazılmıştır; sonuç yalnızca $t\to0$ için geçerlidir.
\end{cevap}

\begin{soru}[Soru 6 \foto{IMG-20261009-WA0019}, \foto{IMG-20261023-WA0046}, \foto{IMG-20261023-WA0037}]
$\mathbf F(t)=\sin t\,\mathbf i+\cos t\,\mathbf j+(1+t)\,\mathbf k$, $\mathbf G(t)=\cos2t\,\mathbf i+\sin3t\,\mathbf j+(e^t+1)\,\mathbf k$ için
$\displaystyle\lim_{t\to0}\mathbf F(t)\cdot\mathbf G(t)$ ve $\displaystyle\lim_{t\to0}\mathbf F(t)\times\mathbf G(t)$'yi bulunuz.
\end{soru}
\begin{cevap}[Çözüm (adım adım)]
\textbf{Adım 1.} $\lim_{t\to0}\mathbf F(t)=(\sin0,\cos0,1+0)=(0,1,1)$; \ $\lim_{t\to0}\mathbf G(t)=(\cos0,\sin0,e^0+1)=(1,0,2)$.

\textbf{Adım 2.} Limit kuralları: çarpımın limiti, limitlerin çarpımıdır.
\[ \lim_{t\to0}\mathbf F\cdot\mathbf G=(0,1,1)\cdot(1,0,2)=0+0+2=2 . \]
\[ \lim_{t\to0}\mathbf F\times\mathbf G=\begin{vmatrix}\mathbf i&\mathbf j&\mathbf k\\0&1&1\\1&0&2\end{vmatrix}=\mathbf i(2-0)-\mathbf j(0-1)+\mathbf k(0-1)=2\,\mathbf i+\mathbf j-\mathbf k . \]
\end{cevap}

\begin{soru}[Soru 7 \foto{IMG-20261009-WA0019}, \foto{IMG-20261009-WA0021}, \foto{IMG-20261023-WA0037}]
$\mathbf F(t)=(|t|+1)\,\mathbf i+t\,[\![t]\!]\,\mathbf j+(2+t^2)\,\mathbf k$ fonksiyonu $t=0$ ve $t=1$ noktalarında sürekli midir?
($[\![t]\!]$: tam değer fonksiyonu.)
\end{soru}
\begin{cevap}[Çözüm (adım adım)]
$\mathbf F$'nin sürekliliği bileşenlerin sürekliliğine denktir. $|t|+1$ ve $2+t^2$ her yerde süreklidir; yalnızca
$g(t)=t\,[\![t]\!]$ incelenir.

\textbf{$t=0$:} $g(0)=0$. Soldan ($-1<t<0$): $[\![t]\!]=-1$, $g(t)=-t\to0$. Sağdan ($0<t<1$): $[\![t]\!]=0$, $g(t)=0$.
Sağ limit $=$ sol limit $=g(0)=0$ $\Rightarrow$ $\mathbf F$, $t=0$'da \textbf{süreklidir}.

\textbf{$t=1$:} Soldan ($0<t<1$): $g(t)=t\cdot0\to0$. Sağdan ($1<t<2$): $g(t)=t\cdot1\to1$. Sol ve sağ limitler farklı, limit yok
$\Rightarrow$ $\mathbf F$, $t=1$'de \textbf{sürekli değildir}.
\end{cevap}

%======================================================================
\subsection{Uzay eğrileri}

\begin{soru}[Soru 8 \foto{IMG-20261009-WA0021}, \foto{IMG-20261023-WA0051}, \foto{IMG-20261023-WA0054}]
$\mathbf r(t)=(2+3t)\,\mathbf i+(4+5t)\,\mathbf j+(6+4t)\,\mathbf k$ eğrisinin cinsi nedir?
\end{soru}
\begin{cevap}
$x=2+3t$, $y=4+5t$, $z=6+4t$; her birinden $t$ çekilip eşitlenirse
\[ \frac{x-2}{3}=\frac{y-4}{5}=\frac{z-6}{4}\ (=t) . \]
Bu, $(2,4,6)$ noktasından geçen ve $\mathbf v=(3,5,4)$ vektörüne paralel \textbf{doğrunun} denklemidir.
\end{cevap}

\begin{soru}[Soru 9 \foto{IMG-20261009-WA0022}, \foto{IMG-20261023-WA0051}, \foto{IMG-20261023-WA0054}]
$\mathbf r(t)=4\cos t\,\mathbf i+4\sin t\,\mathbf j$, $0\le t\le2\pi$ eğrisini tanıyınız.
\end{soru}
\begin{cevap}
$x=4\cos t$, $y=4\sin t$ $\Rightarrow$ $x^2+y^2=16(\cos^2t+\sin^2t)=16$. $xy$ düzleminde, merkezi orijin, yarıçapı $4$ olan \textbf{çember}.
\end{cevap}

\begin{soru}[Soru 10 \foto{IMG-20261009-WA0022}, \foto{IMG-20261023-WA0054}, \foto{IMG-20261023-WA0059}]
$\mathbf r(t)=\cos t\,\mathbf i+\sin t\,\mathbf j+3\,\mathbf k$, $0\le t\le2\pi$ eğrisini tanıyınız.
\end{soru}
\begin{cevap}
$x=\cos t$, $y=\sin t$, $z=3$. $z$ sabit: eğri $z=3$ düzleminde. $x^2+y^2=\cos^2t+\sin^2t=1$ (tahtada ``$y^2+y^2$'' yazılmış,
doğrusu $x^2+y^2$). Sonuç: $z=3$ düzleminde, merkezi $(0,0,3)$, yarıçapı $1$ olan \textbf{çember}.
\end{cevap}

\begin{soru}[Soru 11 \foto{IMG-20261009-WA0025}, \foto{IMG-20261023-WA0054}]
$\mathbf r(t)=2\cos t\,\mathbf i+2\sin t\,\mathbf j+4t\,\mathbf k$ eğrisini inceleyiniz.
\end{soru}
\begin{cevap}[Çözüm (adım adım)]
\textbf{Adım 1.} $x=2\cos t$, $y=2\sin t$, $z=4t$.

\textbf{Adım 2.} $x^2+y^2=4(\cos^2t+\sin^2t)=4$: noktalar $x^2+y^2=4$ silindiri üzerindedir. (Defter ve tahtada
``$4(\cos t+\sin t)$'' yazılmış; kareler unutulmuştur.)

\textbf{Adım 3.} $z=4t$ düzgün artar: $t=1\Rightarrow z=4$, $t=2\Rightarrow z=8$. Bir turda ($t:0\to2\pi$) yükselme $4\cdot2\pi=8\pi$.

\textbf{Sonuç:} Silindire sarılan, adımı $8\pi$ olan \textbf{dairesel helis}.
\end{cevap}

\begin{soru}[Soru 12 \foto{IMG-20261009-WA0025}, \foto{IMG-20261009-WA0020}, \foto{IMG-20261023-WA0039}, \foto{IMG-20261023-WA0059}]
$\mathbf r(t)=(t^2-4)\,\mathbf i+(t^3-4t)\,\mathbf j$, $-2\le t\le2$ eğrisi kapalı mıdır? Basit midir? Orijinde köşe var mıdır?
\end{soru}
\begin{cevap}[Çözüm (adım adım)]
\textbf{Adım 1 (kapalılık).} $\mathbf r(-2)=(0,0)$ ve $\mathbf r(2)=(0,0)$; uçlar çakışır, eğri \textbf{kapalıdır}.

\textbf{Adım 2 (katlı nokta araması).} $\mathbf r(t_1)=\mathbf r(t_2)$ olsun. $x$'lerden $t_1^2=t_2^2$, yani $t_1=t_2$ (aynı nokta) ya da
$t_2=-t_1$. İkinci durumda $y$'ler:
\[ t_1^3-4t_1=-t_1^3+4t_1\ \Rightarrow\ 2t_1^3-8t_1=0\ \Rightarrow\ t_1(t_1^2-4)=0\ \Rightarrow\ t_1\in\{0,\,2,\,-2\}. \]
$t_1=0$ ise $t_2=0$ (aynı parametre); $t_1=\pm2$ ise $t_2=\mp2$ (uç noktalar). Uçlar dışında kesişme yok:
eğri \textbf{basit kapalı eğridir}.

\textbf{Adım 3 (köşe).} $\mathbf r'(t)=2t\,\mathbf i+(3t^2-4)\,\mathbf j$; $\mathbf r'(2)=4\mathbf i+8\mathbf j$, $\mathbf r'(-2)=-4\mathbf i+8\mathbf j$.
Bu vektörler paralel değil (biri diğerinin skaler katı değil): eğri orijinde \textbf{köşe yapar}, düzgün kapanmaz.
\end{cevap}

\begin{soru}[Soru 13 \foto{IMG-20261009-WA0020}, \foto{IMG-20261023-WA0041}, \foto{IMG-20261023-WA0059}]
$\mathbf r(t)=3\cos t\,\mathbf i+3\sin t\,\mathbf j$, $0\le t\le2\pi$ eğrisinin grafiği ve yönü nedir?
\end{soru}
\begin{cevap}
$x^2+y^2=9$: yarıçapı $3$, merkezi orijin olan çember. $t=0:(3,0)$, $t=\pi/2:(0,3)$, $t=\pi:(-3,0)$, $t=3\pi/2:(0,-3)$,
$t=2\pi:(3,0)$. Nokta çemberi \textbf{saat yönünün tersine} dolaşır; bu, pozitif yön kabul edilir.
\end{cevap}

%======================================================================
\subsection{Türev ve türev kuralları}

\begin{soru}[Soru 14 \foto{IMG-20261009-WA0023}, \foto{IMG-20261023-WA0062}]
$\mathbf F(t)=t^2\,\mathbf i+e^{-t}\,\mathbf j+\sin t\,\mathbf k$ için birinci, ikinci ve üçüncü türevleri bulunuz.
\end{soru}
\begin{cevap}
\begin{align*}
\frac{d\mathbf F}{dt}&=2t\,\mathbf i-e^{-t}\,\mathbf j+\cos t\,\mathbf k, &
\frac{d^2\mathbf F}{dt^2}&=2\,\mathbf i+e^{-t}\,\mathbf j-\sin t\,\mathbf k, &
\frac{d^3\mathbf F}{dt^3}&=-e^{-t}\,\mathbf j-\cos t\,\mathbf k .
\end{align*}
(Paydalar $dt^2$, $dt^3$ olmalıdır; tahtada ``$dt$'' yazılmıştır.)
\end{cevap}

\begin{soru}[Soru 15 \foto{IMG-20261009-WA0024}, \foto{IMG-20261023-WA0062}, \foto{IMG-20261023-WA0064}, \foto{IMG-20261023-WA0068}]
$\mathbf r(t)=(1-\cos t)\,\mathbf i+\sin t\,\mathbf j+t\,\mathbf k$ için $\mathbf r'$, $\mathbf r''$ ve $t=\pi/2$'deki değerlerini bulunuz.
\end{soru}
\begin{cevap}
$\mathbf r'(t)=\sin t\,\mathbf i+\cos t\,\mathbf j+\mathbf k$, \ $\mathbf r'(\pi/2)=\mathbf i+0\,\mathbf j+\mathbf k=\mathbf i+\mathbf k$.

$\mathbf r''(t)=\cos t\,\mathbf i-\sin t\,\mathbf j$, \ $\mathbf r''(\pi/2)=-\mathbf j$.

(Tahtada $\mathbf r'(\pi/2)$'de $\mathbf k$ katsayısı ``$0$'' gibi okunuyor; $\mathbf k$ bileşeni sabit $1$'dir.)
\end{cevap}

\begin{soru}[Soru 16 \foto{IMG-20261009-WA0024}, \foto{IMG-20261023-WA0064}, \foto{IMG-20261023-WA0068}]
$\mathbf r(t)=a\cos t\,\mathbf i+a\sin t\,\mathbf j+bt\,\mathbf k$ için $\mathbf r'$, $\mathbf r''$ ve $\mathbf r'\times\mathbf r''$'yü bulunuz.
\end{soru}
\begin{cevap}[Çözüm (adım adım)]
$\mathbf r'(t)=-a\sin t\,\mathbf i+a\cos t\,\mathbf j+b\,\mathbf k$, \quad $\mathbf r''(t)=-a\cos t\,\mathbf i-a\sin t\,\mathbf j$.
\begin{align*}
\mathbf r'\times\mathbf r''&=\begin{vmatrix}\mathbf i&\mathbf j&\mathbf k\\-a\sin t&a\cos t&b\\-a\cos t&-a\sin t&0\end{vmatrix}\\
&=\mathbf i\big(a\cos t\cdot0-b(-a\sin t)\big)-\mathbf j\big(-a\sin t\cdot0-b(-a\cos t)\big)+\mathbf k\big(a^2\sin^2t+a^2\cos^2t\big)\\
&=ab\sin t\,\mathbf i-ab\cos t\,\mathbf j+a^2\,\mathbf k .
\end{align*}
\end{cevap}

\begin{soru}[Soru 17 \foto{IMG-20261009-WA0024}, \foto{IMG-20261023-WA0064}, \foto{IMG-20261023-WA0068}]
$\mathbf u(t)=t\,\mathbf i-t\,\mathbf j+\frac{t^2}{2}\,\mathbf k$, $\mathbf v(t)=\frac{t^4}{4}\,\mathbf i+\frac{t^3}{3}\,\mathbf j+\frac{t^2}{2}\,\mathbf k$ için
$\dfrac{d}{dt}(\mathbf u\times\mathbf v)$'yi bulunuz.
\end{soru}
\begin{cevap}[Çözüm (iki yol)]
\textbf{1. yol.}
\[ \mathbf u\times\mathbf v=\begin{vmatrix}\mathbf i&\mathbf j&\mathbf k\\t&-t&\frac{t^2}{2}\\\frac{t^4}{4}&\frac{t^3}{3}&\frac{t^2}{2}\end{vmatrix}
=\Big(-\frac{t^3}{2}-\frac{t^5}{6}\Big)\mathbf i+\Big(\frac{t^6}{8}-\frac{t^3}{2}\Big)\mathbf j+\Big(\frac{t^4}{3}+\frac{t^5}{4}\Big)\mathbf k . \]
\[ \frac{d}{dt}(\mathbf u\times\mathbf v)=\Big(-\frac{3t^2}{2}-\frac{5t^4}{6}\Big)\mathbf i+\Big(\frac{3t^5}{4}-\frac{3t^2}{2}\Big)\mathbf j+\Big(\frac{4t^3}{3}+\frac{5t^4}{4}\Big)\mathbf k . \]
\textbf{2. yol.} $(\mathbf u\times\mathbf v)'=\mathbf u'\times\mathbf v+\mathbf u\times\mathbf v'$, \ $\mathbf u'=(1,-1,t)$, $\mathbf v'=(t^3,t^2,t)$:
\begin{align*}
\mathbf u'\times\mathbf v&=\Big(-\frac{t^2}{2}-\frac{t^4}{3},\ \frac{t^5}{4}-\frac{t^2}{2},\ \frac{t^3}{3}+\frac{t^4}{4}\Big),\\
\mathbf u\times\mathbf v'&=\Big(-t^2-\frac{t^4}{2},\ \frac{t^5}{2}-t^2,\ t^3+t^4\Big).
\end{align*}
Toplam 1. yoldaki sonucun aynısıdır. \checkmark
\end{cevap}

\begin{dikkat}[Defterdeki hata]
09.10 defterinde (\foto{IMG-20261009-WA0024}) $\mathbf j$ bileşeninin türevi $\frac{3t^5}{2}$ ile yazılmıştır;
$\frac{d}{dt}\frac{t^6}{8}=\frac{3t^5}{4}$'tür. 23.10 tahtasında doğru ($\frac68t^5$) yazılmıştır.
\end{dikkat}

%======================================================================
\subsection{Frenet üçlüsü, hız--ivme ve eğrilik}

\begin{soru}[Soru 18 \foto{IMG-20261015-WA0034}, \foto{IMG-20261023-WA0047}, \foto{IMG-20261023-WA0067}]
$\mathbf r(t)=a\cos t\,\mathbf i+a\sin t\,\mathbf j$ ($a>0$) çemberinin $\mathbf T$, $\mathbf N$, $\mathbf B$ vektörlerini bulunuz.
\end{soru}
\begin{cevap}[Çözüm (adım adım)]
\textbf{Adım 1.} $\mathbf r'(t)=-a\sin t\,\mathbf i+a\cos t\,\mathbf j$, \ $\|\mathbf r'(t)\|=\sqrt{a^2(\sin^2t+\cos^2t)}=a$.

\textbf{Adım 2.} $\mathbf T=\dfrac{\mathbf r'}{\|\mathbf r'\|}=-\sin t\,\mathbf i+\cos t\,\mathbf j$ (defterde işaret unutulmuş: ``$\sin t\,\mathbf i+\cos t\,\mathbf j$'').

\textbf{Adım 3.} $\mathbf T'(t)=-\cos t\,\mathbf i-\sin t\,\mathbf j$, \ $\|\mathbf T'\|=1$ $\Rightarrow$ $\mathbf N=-\cos t\,\mathbf i-\sin t\,\mathbf j=-\frac1a\mathbf r(t)$ (merkeze doğru).

\textbf{Adım 4.}
\[ \mathbf B=\mathbf T\times\mathbf N=\begin{vmatrix}\mathbf i&\mathbf j&\mathbf k\\-\sin t&\cos t&0\\-\cos t&-\sin t&0\end{vmatrix}=(\sin^2t+\cos^2t)\,\mathbf k=\mathbf k . \]
\end{cevap}

\begin{soru}[Soru 19 \foto{IMG-20261015-WA0036}, \foto{IMG-20261023-WA0049}, \foto{IMG-20261023-WA0067}]
$a>0$, $\mathbf r(t)=a\cos t\,\mathbf i+a\sin t\,\mathbf j$ olan hareketlinin hız vektörü konum vektörüne dik midir?
\end{soru}
\begin{cevap}
$\mathbf v=\mathbf r'=-a\sin t\,\mathbf i+a\cos t\,\mathbf j$; \ $\mathbf v\cdot\mathbf r=-a^2\sin t\cos t+a^2\sin t\cos t=0$.
Her $t$ için $\mathbf v\perp\mathbf r$ (çemberde hız yarıçapa diktir).
\end{cevap}

\begin{soru}[Soru 20 \foto{IMG-20261015-WA0036}, \foto{IMG-20261023-WA0049}, \foto{IMG-20261023-WA0067}]
$\mathbf r(t)=5\cos t\,\mathbf i+5\sin t\,\mathbf j+6t^2\,\mathbf k$ olan hareketlinin hızı ve ivmesi nedir? Hız ve ivme ne zaman diktir?
\end{soru}
\begin{cevap}
$\mathbf v(t)=-5\sin t\,\mathbf i+5\cos t\,\mathbf j+12t\,\mathbf k$, \ $\mathbf a(t)=-5\cos t\,\mathbf i-5\sin t\,\mathbf j+12\,\mathbf k$.
\[ \mathbf a\cdot\mathbf v=25\sin t\cos t-25\sin t\cos t+144t=144t=0\ \Rightarrow\ t=0 . \]
\end{cevap}

\begin{soru}[Soru 21 \foto{IMG-20261015-WA0036}, \foto{IMG-20261023-WA0049}, \foto{IMG-20261023-WA0038}]
Parametrik denklemleri $x=e^{-t}$, $y=2\cos3t$, $z=2\sin3t$ olan eğri boyunca bir cisim hareket ediyor ($t$: zaman).
(a) Hız ve ivme? (b) $t=0$ iken büyüklükleri? (c) $t=0$ iken birim teğet vektör?
\end{soru}
\begin{cevap}[Çözüm (adım adım)]
$\mathbf r(t)=e^{-t}\,\mathbf i+2\cos3t\,\mathbf j+2\sin3t\,\mathbf k$.

\textbf{(a)} Zincir kuralıyla ($\frac{d}{dt}\cos3t=-3\sin3t$):
$\mathbf v=-e^{-t}\,\mathbf i-6\sin3t\,\mathbf j+6\cos3t\,\mathbf k$, \ $\mathbf a=e^{-t}\,\mathbf i-18\cos3t\,\mathbf j-18\sin3t\,\mathbf k$.

\textbf{(b)} $\mathbf v(0)=-\mathbf i+6\mathbf k$, $\|\mathbf v(0)\|=\sqrt{37}$; \ $\mathbf a(0)=\mathbf i-18\mathbf j$, $\|\mathbf a(0)\|=\sqrt{325}=5\sqrt{13}$.

\textbf{(c)} $\mathbf T(0)=\dfrac{\mathbf v(0)}{\|\mathbf v(0)\|}=-\dfrac{1}{\sqrt{37}}\,\mathbf i+\dfrac{6}{\sqrt{37}}\,\mathbf k$.
\end{cevap}

\begin{soru}[Soru 22 \foto{IMG-20261015-WA0035}, \foto{IMG-20261023-WA0052}]
Bir nokta $\mathbf r=\mathbf r(t)$ eğrisi boyunca hareket ediyor; $s$ alınan yol, $\upsilon=\|d\mathbf r/dt\|=ds/dt$.
$\mathbf a=\dfrac{d\upsilon}{dt}\mathbf T+\dfrac{\upsilon^2}{\rho}\mathbf N$ olduğunu gösteriniz.
\end{soru}
\begin{cevap}[Çözüm (adım adım)]
\textbf{Adım 1.} $\mathbf v=\dfrac{d\mathbf r}{dt}=\dfrac{d\mathbf r}{ds}\dfrac{ds}{dt}=\upsilon\,\mathbf T$.

\textbf{Adım 2.} Çarpım kuralı: $\mathbf a=\dfrac{d\mathbf v}{dt}=\dfrac{d\upsilon}{dt}\mathbf T+\upsilon\dfrac{d\mathbf T}{dt}$.

\textbf{Adım 3.} Zincir kuralı: $\dfrac{d\mathbf T}{dt}=\dfrac{d\mathbf T}{ds}\dfrac{ds}{dt}=\upsilon\dfrac{d\mathbf T}{ds}$; eğrilik tanımından $\dfrac{d\mathbf T}{ds}=\kappa\mathbf N=\dfrac1\rho\mathbf N$.

\textbf{Adım 4.} Yerine yazılırsa $\mathbf a=\dfrac{d\upsilon}{dt}\mathbf T+\upsilon^2\dfrac{d\mathbf T}{ds}=\dfrac{d\upsilon}{dt}\mathbf T+\dfrac{\upsilon^2}{\rho}\mathbf N$. \hfill$\blacksquare$
\end{cevap}

\begin{soru}[Soru 23 \foto{IMG-20261015-WA0038}, \foto{IMG-20261023-WA0052}]
$\mathbf r(t)=a\cos t\,\mathbf i+a\sin t\,\mathbf j$ çemberinin eğriliğini ve eğrilik yarıçapını bulunuz; $t$ zaman ise ivmeyi
$\mathbf T$, $\mathbf N$ bileşenlerine ayırınız.
\end{soru}
\begin{cevap}[Çözüm (adım adım)]
Soru 18'den $\mathbf T'(t)=-\cos t\,\mathbf i-\sin t\,\mathbf j$, $\frac{ds}{dt}=a$.
\[ \frac{d\mathbf T}{ds}=\frac{d\mathbf T/dt}{ds/dt}=\frac1a(-\cos t\,\mathbf i-\sin t\,\mathbf j)\ \Rightarrow\ \kappa=\frac1a,\quad\rho=a . \]
Sürat $\upsilon=a$ sabit $\Rightarrow\frac{d\upsilon}{dt}=0$. \
$\mathbf a(t)=\mathbf r''(t)=-a\cos t\,\mathbf i-a\sin t\,\mathbf j=a\,\mathbf N=\dfrac{\upsilon^2}{\rho}\mathbf N$ ($a^2/a=a$ \checkmark).
\end{cevap}

%======================================================================
\subsection{Düzgün eğri, integral ve yay uzunluğu}

\begin{soru}[Soru 24 \foto{IMG-20261015-WA0038}, \foto{IMG-20261023-WA0053}]
$\mathbf r(t)=a(t-\sin t)\,\mathbf i+a(1-\cos t)\,\mathbf j$, $-2\pi\le t\le2\pi$ eğrisi düzgün müdür, parçalı düzgün müdür?
\end{soru}
\begin{cevap}[Çözüm (adım adım)]
\textbf{Adım 1.} $\mathbf r'(t)=a(1-\cos t)\,\mathbf i+a\sin t\,\mathbf j$ süreklidir.

\textbf{Adım 2.} $\mathbf r'=\mathbf 0\iff\cos t=1$ ve $\sin t=0\iff t=2k\pi$. Aralıkta: $t=-2\pi,\,0,\,2\pi$.

\textbf{Adım 3.} $t=0$ iç nokta olduğundan eğri $[-2\pi,2\pi]$'de düzgün değildir. Ama $[-2\pi,0]$ ve $[0,2\pi]$ alt aralıklarının
iç noktalarında $\mathbf r'\neq\mathbf 0$: her biri düzgündür. Eğri \textbf{parçalı düzgündür}; $x=-2\pi a,0,2\pi a$ sivri uçlardır.
\end{cevap}

\begin{soru}[Soru 25 \foto{IMG-20261015-WA0037}, \foto{IMG-20261023-WA0055}]
$\mathbf r(t)=2t\,\mathbf i+e^t\,\mathbf j+4\sin2t\,\mathbf k$ fonksiyonunun belirsiz integralini bulunuz.
\end{soru}
\begin{cevap}
\[ \int\mathbf r\,dt=t^2\,\mathbf i+e^t\,\mathbf j-2\cos2t\,\mathbf k+\mathbf C=(t^2+c_1)\,\mathbf i+(e^t+c_2)\,\mathbf j+(-2\cos2t+c_3)\,\mathbf k . \]
($4\int\sin2t\,dt=4\cdot\big(-\frac12\cos2t\big)=-2\cos2t$.)
\end{cevap}

\begin{soru}[Soru 26 \foto{IMG-20261015-WA0037}, \foto{IMG-20261023-WA0055}]
$\mathbf r'(t)=2t\,\mathbf i+\sin t\,\mathbf j+e^t\,\mathbf k$ ve $\mathbf r(0)=2\,\mathbf i+2\,\mathbf j$ koşullarını sağlayan $\mathbf r(t)$'yi bulunuz.
\end{soru}
\begin{cevap}[Çözüm (adım adım)]
\textbf{Adım 1.} $\mathbf r(t)=t^2\,\mathbf i-\cos t\,\mathbf j+e^t\,\mathbf k+\mathbf C$.

\textbf{Adım 2.} $t=0$: $-\mathbf j+\mathbf k+\mathbf C=2\mathbf i+2\mathbf j$ $\Rightarrow$ $\mathbf C=2\mathbf i+3\mathbf j-\mathbf k$.

\textbf{Adım 3.} $\boxed{\mathbf r(t)=(t^2+2)\,\mathbf i+(3-\cos t)\,\mathbf j+(e^t-1)\,\mathbf k}$. \
Kontrol: $\mathbf r(0)=2\mathbf i+2\mathbf j+0\mathbf k$ \checkmark
\end{cevap}

\begin{soru}[Soru 27 \foto{IMG-20261015-WA0039}, \foto{IMG-20261023-WA0040}]
$\mathbf F(t)=2t\,\mathbf i+(1+t^2)\,\mathbf j+\sin\pi t\,\mathbf k$ için $\displaystyle\int_0^1\mathbf F(t)\,dt$ nedir?
\end{soru}
\begin{cevap}
\[ \int_0^1\mathbf F\,dt=t^2\Big|_0^1\mathbf i+\Big(t+\frac{t^3}{3}\Big)\Big|_0^1\mathbf j-\frac{\cos\pi t}{\pi}\Big|_0^1\mathbf k
=\mathbf i+\frac43\,\mathbf j-\frac{-1-1}{\pi}\,\mathbf k=\mathbf i+\frac43\,\mathbf j+\frac2\pi\,\mathbf k . \]
\end{cevap}

\begin{soru}[Soru 28 \foto{IMG-20261015-WA0039}, \foto{IMG-20261023-WA0040}, \foto{IMG-20261023-WA0058}]
$\mathbf r(t)=5\cos t\,\mathbf i+5\sin t\,\mathbf j+12t\,\mathbf k$, $0\le t\le2\pi$ helis parçasının uzunluğu nedir?
\end{soru}
\begin{cevap}[Çözüm (adım adım)]
\textbf{Adım 1.} $\mathbf r'(t)=-5\sin t\,\mathbf i+5\cos t\,\mathbf j+12\,\mathbf k$.

\textbf{Adım 2.} $\|\mathbf r'(t)\|=\sqrt{25\sin^2t+25\cos^2t+144}=\sqrt{169}=13$.

\textbf{Adım 3.} $\ell=\int_0^{2\pi}13\,dt=26\pi$.

\textbf{Sağlama (tahtadaki gibi):} Silindir açılınca helis, kenarları $2\pi\cdot5=10\pi$ ve $12\cdot2\pi=24\pi$ olan dikdörtgenin
köşegenidir: $\sqrt{(10\pi)^2+(24\pi)^2}=26\pi$ \checkmark
\end{cevap}

%======================================================================
\subsection{Çalışma soruları}

Aşağıdaki problemler 15.10 defterinde ve 23.10 tahtasında (\foto{IMG-20261015-WA0039}, \foto{IMG-20261015-WA0040},
\foto{IMG-20261023-WA0058}) çözümsüz olarak verilmiştir.

\begin{soru}[Problem 1 (derste çalışma sorusu olarak verildi) \foto{IMG-20261015-WA0039}, \foto{IMG-20261023-WA0058}]
Türevlerini bulunuz:
(a) $\mathbf F(t)=t\,\mathbf i+t^2\,\mathbf j+t^3\,\mathbf k$; \
(b) $\mathbf F(t)=(t\ln t-t)\,\mathbf i+e^t\,\mathbf j+e^t\,\mathbf k$; \
(c) $\mathbf F(t)=\arcsin(4t)\,\mathbf i-3\arctan(2t-1)\,\mathbf j+7\ln(3t^2)\,\mathbf k$.
\end{soru}
\begin{cevap}[Çözüm (adım adım)]
\textbf{(a)} $\mathbf F'(t)=\mathbf i+2t\,\mathbf j+3t^2\,\mathbf k$.

\textbf{(b)} $\frac{d}{dt}(t\ln t-t)=\ln t+t\cdot\frac1t-1=\ln t$ (çarpım kuralı). Böylece $\mathbf F'(t)=\ln t\,\mathbf i+e^t\,\mathbf j+e^t\,\mathbf k$ \ ($t>0$).

\textbf{(c)} $\frac{d}{dt}\arcsin u=\frac{u'}{\sqrt{1-u^2}}$, $\frac{d}{dt}\arctan u=\frac{u'}{1+u^2}$, $\frac{d}{dt}\ln u=\frac{u'}{u}$:
\[ \mathbf F'(t)=\frac{4}{\sqrt{1-16t^2}}\,\mathbf i-\frac{6}{1+(2t-1)^2}\,\mathbf j+\frac{14}{t}\,\mathbf k\qquad\Big(\tfrac{7\cdot6t}{3t^2}=\tfrac{14}{t}\Big). \]
\textit{Not:} 15.10 defterinde (b)'nin $\mathbf k$ bileşenindeki üs belirsizdir (``$e^{-t}$'' de olabilir); 23.10 tahtasında açıkça $e^t$
okunmaktadır. $e^{-t}$ olsaydı türevi $-e^{-t}\,\mathbf k$ olurdu.
\end{cevap}

\begin{soru}[Problem 2 (derste çalışma sorusu olarak verildi) \foto{IMG-20261015-WA0039}, \foto{IMG-20261023-WA0058}]
$\mathbf F$ iki kez türevlenebilen bir fonksiyon olsun. $[\mathbf F(t)\times\mathbf F'(t)]'=\mathbf F(t)\times\mathbf F''(t)$ olduğunu gösteriniz.
\end{soru}
\begin{cevap}
Vektörel çarpımın türev kuralıyla
\[ [\mathbf F\times\mathbf F']'=\mathbf F'\times\mathbf F'+\mathbf F\times\mathbf F'' . \]
Bir vektörün kendisiyle vektörel çarpımı sıfırdır ($\sin0=0$): $\mathbf F'\times\mathbf F'=\mathbf 0$. Dolayısıyla
$[\mathbf F\times\mathbf F']'=\mathbf F\times\mathbf F''$. \hfill$\blacksquare$
\end{cevap}

\begin{soru}[Problem 3 (derste çalışma sorusu olarak verildi) \foto{IMG-20261015-WA0040}, \foto{IMG-20261023-WA0058}]
$\mathbf r(t)=e^t\,\mathbf i+e^{-t}\,\mathbf j+\sqrt2\,t\,\mathbf k$ eğrisinin $t=0$ noktasındaki Frenet üçlüsünü hesaplayınız.
\end{soru}
\begin{cevap}[Çözüm (adım adım)]
\textbf{Adım 1 ($\mathbf T$).} $\mathbf r'(t)=(e^t,-e^{-t},\sqrt2)$, $\mathbf r'(0)=(1,-1,\sqrt2)$, $\|\mathbf r'(0)\|=\sqrt{1+1+2}=2$.
\[ \mathbf T(0)=\Big(\frac12,\,-\frac12,\,\frac{\sqrt2}{2}\Big). \]
\textbf{Adım 2 ($\mathbf B$).} $\mathbf r''(t)=(e^t,e^{-t},0)$, $\mathbf r''(0)=(1,1,0)$.
\[ \mathbf r'\times\mathbf r''=\begin{vmatrix}\mathbf i&\mathbf j&\mathbf k\\1&-1&\sqrt2\\1&1&0\end{vmatrix}=(0-\sqrt2)\,\mathbf i-(0-\sqrt2)\,\mathbf j+(1+1)\,\mathbf k=(-\sqrt2,\sqrt2,2), \]
normu $\sqrt{2+2+4}=2\sqrt2$. \quad $\mathbf B(0)=\Big(-\dfrac12,\,\dfrac12,\,\dfrac{\sqrt2}{2}\Big)$.

\textbf{Adım 3 ($\mathbf N$).} $\mathbf N=\mathbf B\times\mathbf T$:
\[ \mathbf N(0)=\begin{vmatrix}\mathbf i&\mathbf j&\mathbf k\\-\frac12&\frac12&\frac{\sqrt2}2\\\frac12&-\frac12&\frac{\sqrt2}2\end{vmatrix}
=\Big(\frac{\sqrt2}{4}+\frac{\sqrt2}{4}\Big)\mathbf i-\Big(-\frac{\sqrt2}{4}-\frac{\sqrt2}{4}\Big)\mathbf j+\Big(\frac14-\frac14\Big)\mathbf k=\Big(\frac{\sqrt2}{2},\frac{\sqrt2}{2},0\Big). \]
\textbf{Sağlama (tanımla).} $\|\mathbf r'(t)\|=\sqrt{e^{2t}+e^{-2t}+2}=e^t+e^{-t}$; bu fonksiyonun $t=0$'daki türevi $0$ olduğundan
$\mathbf T'(0)=\frac{\mathbf r''(0)}{\|\mathbf r'(0)\|}=\frac12(1,1,0)$ ve $\mathbf N(0)=\frac{(1,1,0)}{\sqrt2}$ \checkmark. Ek olarak
$\kappa(0)=\frac{2\sqrt2}{2^3}=\frac{\sqrt2}{4}$.
\end{cevap}

\begin{soru}[Problem 4 (derste çalışma sorusu olarak verildi) \foto{IMG-20261015-WA0040}, \foto{IMG-20261023-WA0058}]
İntegralleri hesaplayınız:
(a) $\displaystyle\int_0^1\big[4t^3\,\mathbf i+3\,\mathbf j+2(t+1)\,\mathbf k\big]dt$; \
(b) $\displaystyle\int_1^2\Big[(6-6t)\,\mathbf i+3\sqrt t\,\mathbf j+\frac4{t^2}\,\mathbf k\Big]dt$; \
(c) $\displaystyle\int_{-\pi/4}^{\pi/4}\big[\sin t\,\mathbf i+(1+\cos t)\,\mathbf j+\sec^2t\,\mathbf k\big]dt$.
\end{soru}
\begin{cevap}[Çözüm (adım adım)]
\textbf{(a)} $\int_0^14t^3dt=1$; $\int_0^13\,dt=3$; $\int_0^12(t+1)dt=[t^2+2t]_0^1=3$. \ Sonuç: $\mathbf i+3\,\mathbf j+3\,\mathbf k$.

\textbf{(b)} $\int_1^2(6-6t)dt=[6t-3t^2]_1^2=0-3=-3$; \ $\int_1^23\sqrt t\,dt=[2t^{3/2}]_1^2=4\sqrt2-2$; \
$\int_1^2\frac4{t^2}dt=\big[-\frac4t\big]_1^2=-2+4=2$. \ Sonuç: $-3\,\mathbf i+(4\sqrt2-2)\,\mathbf j+2\,\mathbf k$.

\textbf{(c)} $\sin t$ tek fonksiyon, simetrik aralıkta integrali $0$; \
$\int_{-\pi/4}^{\pi/4}(1+\cos t)dt=\frac\pi2+[\sin t]_{-\pi/4}^{\pi/4}=\frac\pi2+\sqrt2$; \
$\int_{-\pi/4}^{\pi/4}\sec^2t\,dt=[\tan t]_{-\pi/4}^{\pi/4}=2$. \ Sonuç: $\big(\frac\pi2+\sqrt2\big)\mathbf j+2\,\mathbf k$.
\end{cevap}

%======================================================================
\subsection{Ek alıştırmalar}

\begin{soru}[\ekal Ek alıştırma 1]
$\mathbf F(t)=\sqrt{t-1}\,\mathbf i+\dfrac{1}{t-3}\,\mathbf j+\ln(5-t)\,\mathbf k$ fonksiyonunun tanım kümesini bulunuz.
\end{soru}
\begin{cevap}
$T_f=[1,\infty)$, $T_g=\mathbb R\setminus\{3\}$, $T_h=(-\infty,5)$. Kesişim: $T_{\mathbf F}=[1,3)\cup(3,5)$.
\end{cevap}

\begin{soru}[\ekal Ek alıştırma 2]
$\displaystyle\lim_{t\to0}\Big[\frac{1-\cos t}{t^2}\,\mathbf i+\frac{e^{2t}-1}{t}\,\mathbf j+\frac{\tan t}{t}\,\mathbf k\Big]$ limitini bulunuz.
\end{soru}
\begin{cevap}
Her bileşen $\frac00$ belirsizliğidir; L'Hospital (ya da bilinen limitler):
$\frac{1-\cos t}{t^2}\to\frac{\sin t}{2t}\to\frac12$; \ $\frac{e^{2t}-1}{t}\to\frac{2e^{2t}}{1}\to2$; \ $\frac{\tan t}{t}\to\frac{\sec^2t}{1}\to1$.
Limit: $\frac12\,\mathbf i+2\,\mathbf j+\mathbf k$.
\end{cevap}

\begin{soru}[\ekal Ek alıştırma 3]
$\mathbf r''(t)=6t\,\mathbf i+2\,\mathbf j+e^t\,\mathbf k$, $\mathbf r'(0)=\mathbf j+\mathbf k$, $\mathbf r(0)=\mathbf i$ ise $\mathbf r(t)$'yi bulunuz.
\end{soru}
\begin{cevap}
$\mathbf r'=3t^2\,\mathbf i+2t\,\mathbf j+e^t\,\mathbf k+\mathbf C_1$; $\mathbf r'(0)=\mathbf k+\mathbf C_1=\mathbf j+\mathbf k\Rightarrow\mathbf C_1=\mathbf j$:
$\mathbf r'=3t^2\,\mathbf i+(2t+1)\,\mathbf j+e^t\,\mathbf k$.

$\mathbf r=t^3\,\mathbf i+(t^2+t)\,\mathbf j+e^t\,\mathbf k+\mathbf C_2$; $\mathbf r(0)=\mathbf k+\mathbf C_2=\mathbf i\Rightarrow\mathbf C_2=\mathbf i-\mathbf k$:
\[ \mathbf r(t)=(t^3+1)\,\mathbf i+(t^2+t)\,\mathbf j+(e^t-1)\,\mathbf k . \]
\end{cevap}

\begin{soru}[\ekal Ek alıştırma 4]
$\mathbf r(t)=t^3\,\mathbf i+3t^2\,\mathbf j+6t\,\mathbf k$, $0\le t\le1$ eğrisinin uzunluğunu bulunuz.
\end{soru}
\begin{cevap}
$\mathbf r'=(3t^2,6t,6)$, $\|\mathbf r'\|^2=9t^4+36t^2+36=9(t^2+2)^2$, $\|\mathbf r'\|=3(t^2+2)$.
\[ \ell=\int_0^13(t^2+2)\,dt=3\Big(\frac13+2\Big)=7 . \]
\end{cevap}

\begin{soru}[\ekal Ek alıştırma 5]
$\mathbf r(t)=3\cos t\,\mathbf i+3\sin t\,\mathbf j+4t\,\mathbf k$ helisinin $\mathbf T$, $\mathbf N$, $\mathbf B$ vektörlerini, eğriliğini ve burulmasını bulunuz.
\end{soru}
\begin{cevap}[Çözüm (adım adım)]
$\mathbf r'=(-3\sin t,3\cos t,4)$, $\|\mathbf r'\|=5$: \ $\mathbf T=\frac15(-3\sin t,3\cos t,4)$.

$\mathbf T'=\frac15(-3\cos t,-3\sin t,0)$, $\|\mathbf T'\|=\frac35$: \ $\mathbf N=(-\cos t,-\sin t,0)$.

$\mathbf B=\mathbf T\times\mathbf N=\frac15\begin{vmatrix}\mathbf i&\mathbf j&\mathbf k\\-3\sin t&3\cos t&4\\-\cos t&-\sin t&0\end{vmatrix}=\frac15(4\sin t,-4\cos t,3)$.

$\kappa=\dfrac{\|\mathbf T'\|}{\|\mathbf r'\|}=\dfrac{3/5}{5}=\dfrac{3}{25}$ \ (helis formülü $\frac{a}{a^2+b^2}=\frac{3}{25}$ \checkmark). \
$\tau=\dfrac{b}{a^2+b^2}=\dfrac{4}{25}$ \ (kontrol: $(\mathbf r'\times\mathbf r'')\cdot\mathbf r'''=a^2b=36$, $\|\mathbf r'\times\mathbf r''\|^2=a^2(a^2+b^2)=225$, $36/225=4/25$ \checkmark).
\end{cevap}

\begin{soru}[\ekal Ek alıştırma 6]
$\mathbf r(t)=t\,\mathbf i+2t\,\mathbf j+t^2\,\mathbf k$ boyunca hareket eden cismin ivmesinin teğet ve normal bileşenlerini bulunuz.
\end{soru}
\begin{cevap}
$\mathbf v=(1,2,2t)$, $\mathbf a=(0,0,2)$, $\|\mathbf v\|=\sqrt{5+4t^2}$.
\[ a_T=\frac{\mathbf v\cdot\mathbf a}{\|\mathbf v\|}=\frac{4t}{\sqrt{5+4t^2}},\qquad
\mathbf v\times\mathbf a=(4,-2,0),\quad a_N=\frac{\|\mathbf v\times\mathbf a\|}{\|\mathbf v\|}=\frac{2\sqrt5}{\sqrt{5+4t^2}} . \]
Kontrol: $a_T^2+a_N^2=\frac{16t^2+20}{5+4t^2}=4=\|\mathbf a\|^2$ \checkmark
\end{cevap}

\begin{soru}[\ekal Ek alıştırma 7]
$y=\ln x$ eğrisinin eğriliğini bulunuz. Eğrilik hangi noktada en büyüktür?
\end{soru}
\begin{cevap}
$f'=\frac1x$, $f''=-\frac1{x^2}$:
\[ \kappa(x)=\frac{1/x^2}{(1+1/x^2)^{3/2}}=\frac{x}{(x^2+1)^{3/2}}\quad(x>0). \]
$\kappa'(x)=\dfrac{(x^2+1)-3x^2}{(x^2+1)^{5/2}}=\dfrac{1-2x^2}{(x^2+1)^{5/2}}=0\Rightarrow x=\dfrac1{\sqrt2}$ (işaret $+$'dan $-$'ye döner: maksimum).
En büyük eğrilik $\big(\frac1{\sqrt2},-\frac12\ln2\big)$ noktasındadır: $\kappa_{\max}=\dfrac{1/\sqrt2}{(3/2)^{3/2}}=\dfrac{2\sqrt3}{9}$.
\end{cevap}

\begin{soru}[\ekal Ek alıştırma 8]
Bir top yerden $v_0=20$ m/s süratle, yatayla $60^\circ$ açıyla atılıyor ($g\approx10$ m/s$^2$). Konum vektörünü, uçuş süresini, menzili ve
en büyük yüksekliği bulunuz.
\end{soru}
\begin{cevap}
$\mathbf a=-10\,\mathbf j$, $\mathbf v(0)=20\cos60^\circ\,\mathbf i+20\sin60^\circ\,\mathbf j=10\,\mathbf i+10\sqrt3\,\mathbf j$, $\mathbf r(0)=\mathbf 0$.
İki kez integral: $\mathbf r(t)=10t\,\mathbf i+(10\sqrt3\,t-5t^2)\,\mathbf j$.
\begin{itemize}
\item Uçuş süresi: $10\sqrt3\,t-5t^2=0\Rightarrow t_u=2\sqrt3\approx3{,}46$ s.
\item Menzil: $x(2\sqrt3)=20\sqrt3\approx34{,}6$ m \ ($\frac{v_0^2\sin120^\circ}{g}=\frac{400\cdot\sqrt3/2}{10}$ \checkmark).
\item En büyük yükseklik: $y'=10\sqrt3-10t=0\Rightarrow t=\sqrt3$, $H=30-15=15$ m.
\end{itemize}
\end{cevap}
